\documentclass[11pt]{letter}
\usepackage[T1]{fontenc}
\usepackage[utf8]{inputenc}
\usepackage[margin=1in]{geometry}
\usepackage{lmodern}
\usepackage{hyperref}
\hypersetup{hidelinks}
\setlength{\emergencystretch}{2em}

\signature{Mingyuan Xie\\Zhengxun Wu\\
Qingdao University of Technology\\
Corresponding authors\\
\href{mailto:jhj208436@gmail.com}{jhj208436@gmail.com}\\
\href{mailto:2276815088@qq.com}{2276815088@qq.com}}
\address{Qingdao University of Technology\\Qingdao, Shandong, China}
\date{11 August 2026}

\begin{document}
\begin{letter}{Editor-in-Chief\\\emph{Concurrency and Computation: Practice and Experience}}

\opening{Dear Editor-in-Chief,}

We submit the manuscript entitled ``Concurrency-Aware Permission Decisions in
Multi-Tenant LLM Agent Runtimes'' for consideration as a research article in
\emph{Concurrency and Computation: Practice and Experience}.

The manuscript studies a concurrency failure in a multi-tenant agent runtime:
tenant work can fill a bounded executor and delay a permission decision until
its deadline, while a permissive timeout fallback converts executor
unavailability into unsafe admission. We present Neuron, a Java~21 runtime
architecture that reserves permission capacity and exports a resource-pressure
signal to the escalation policy. The paper's primary contribution is a
systems characterization of executor contention, queueing, deadline semantics,
and permission-verdict behavior; the security consequence is the measured
failure mode rather than a claim of a new authorization theory.

The evaluation compares shared fail-open, shared fail-closed, isolated, and
coupled architectures. It contains five independent JVM runs on two recorded
environments, deadline and attacker-pressure surfaces, capacity sweeps, mixed
load QoS, fixed API-trace replay, a process-separated Java permission service,
and an additional executor-focused Windows replication. Raw JSONL observations,
environment records, analysis scripts, generated values, and reproduction
instructions are publicly archived at
\url{https://github.com/dadawer-web/ouisani/tree/ccpe-repro-20260811/ccpe-repro-20260811}.
The process-separated service is explicitly a same-host loopback experiment;
we do not claim WAN, multi-node, GPU, or live-production behavior.

We believe the manuscript fits the journal's emphasis on practical concurrent
and distributed systems because it reports an implemented Java runtime,
executor-level measurements, architecture baselines, and reproducible
engineering trade-offs. The manuscript is original, is not under consideration
elsewhere, and has been approved by both authors.

Thank you for your consideration.

\closing{Sincerely,}
\end{letter}
\end{document}
